% =================================================================
\section{2.1 선형회귀: 모델, 비용함수, 경사하강법}
\begin{frame}{첫 완전한 알고리즘: 세 부분 뼈대}
  지난 장에서 세운 것은 알고리즘이 아니라 \textbf{문제의 정의}였다.
  이번 장은 1.2의 예고(``같은 선형모델을 숫자로 읽으면 회귀'')를
  \textbf{이 학기의 첫 완전한 알고리즘}으로 실행하는 장이다.
  \vskip1em
  알고리즘의 뼈대는 항상 같은 세 부류:
  \begin{enumerate}
    \item \kb{모델} --- 입력을 출력으로 매핑하는 함수의 \emph{형태}
    \item \kb{비용함수}(loss) --- ``얼마나 틀렸는지''를 숫자로 재는 재
    \item \kb{최적화 알고리즘} --- 파라미터를 그 숫자가 줄어드는
      방향으로 바꾸는 방법
  \end{enumerate}
  \vskip0.6em
  셋이 모여야 비로소 ``알고리즘''. 이 틀은 \textbf{학기 내내 이름만 바꿔
  반복}된다. 그 사이사이에는 손으로, 또는 몇 줄의 코드로 재현 가능한
  \textbf{작은 수치 실험}이 들어간다.
\end{frame}
